\begin{abstract}

    Blockchain scalability is one of the issues that concerns its current
    adopters. The current popular blockchains have initially been designed with
    imperfections that introduce fundamental bottlenecks which limit their
    ability to have a higher throughput and a lower latency.

    One of the major bottlenecks for existing blockchain technologies is fast block
    propagation. A faster block propagation enables a miner to reach a majority
    of the network within a time constraint and therefore leading to a lower
    orphan rate and better profitability. In order to attain a throughput that
    could compete with the current state of the art transaction processing,
    while also keeping the block intervals same as today, a 24.3 Gigabyte
    block will be required every 10 minutes with an average transaction size of
    500 bytes, which translates to 48600000 transactions every 10 minutes or about 81000 transactions per second.

    In order to synchronize such large blocks faster across the network while maintaining 
    consensus by keeping the orphan rate below 50\%, the thesis
    proposes to aggregate partial block data from multiple nodes using digital
    fountain codes. The advantages of using a fountain code is that all connected peers can send
    part of data in an encoded form. When the receiving peer has enough data, it
    then decodes the information to reconstruct the block. Along with them
    sending only part information, the data can be relayed over UDP, instead of
    TCP, improving upon the speed of propagation  in the current blockchains.
    Fountain codes applied in this research are Raptor codes, which allow
    construction of infinite decoding symbols. The research, when applied to
    blockchains, increases success rate of block delivery on decode failures.
      
\end{abstract}
